% One Britannica scan as extracted by Nemotron Parse 3.5. Filled by src/nemotron_challenge/stages/latex.py.
% The TITLE, BODY and SOURCE markers (each wrapped in two percent signs) are replaced. Compile with XeTeX (tectonic).
\documentclass[10pt]{article}
\usepackage[a4paper,margin=18mm]{geometry}
\usepackage{fontspec}
\usepackage{multicol}
\usepackage{graphicx}
\usepackage[hidelinks]{hyperref}
\setlength{\parindent}{0pt}
\setlength{\parskip}{4pt}
\setlength{\columnsep}{8mm}
\begin{document}
\section*{Historical source: BLOW-PIPE}
\begin{multicols}{2}
\subsection*{\textit{continued:} BLOW-PIPE}
to the axis of the blow-pipe. A consists of two pieces, one of which fits air-tight into the other, and may be turned out as six, so that the pipe of issue may be made to form different angles with the axis of the blow-pipe, the position of the matter under exponent may require. Name consists of watering in a suite of incidences. Many substances, both of the vegetable, animals, and mineral kingdoms, have the quality of giving out this inconcedes vapour. For domestic uses, and for the arms, or granized bodies of the vegetable and animal kingdom along are employed to produce it; such as oils, some of which are sold, others think, at the most temperature of its skin, alcohol, other wood, and pricool. This latter, though found among minerals is composed of organized matter channel and rendered bitumen by a particular process of decomposition. The blow-pipe, by directing the manifestances particles of which the flame consistas so as a strike against and surround a small body, produces the effect of heating the body considerably. The flame used with the blow-pipe may either the the flame of an oil or spirit, or of the smaller, the flames of the carbonated hydrogen gas proceeding from the distillation of the pitcoal, is also found advantageous for this purpose. To order to use the blow-pipe, the breath impelled through it is to be directed across the flame of a lamp or quantity, applying this orifice from which the dirt issues a little above on the upper end of the neck; and a jet of flame is thus formed as represented at fig. 8. This jet is made to fall on the body to be heated. The operator may be continued for a considerable length than time; an minister reputed blast: a kept up by the muscular action of the bloods, whilst the ordinary separation goes on through the nose; and a little practice is sufficient to enable the opmore to succeed. The jet of flame is conodal, internal, blue, and externally yellow. By more or less immersion in this jet of flame the subject of operation receives a greater or less degree of feet, but becomes evident in a greater or less degree. If a head of bar, containing oxide of manganese, be kept fused for some time in the inner flame, the bead become colorant, when it is all forward kept fused in the outer flame, the manganese ecquires more oxygen, and the bead becomes of a violet color. This violet colour may be made to appear more speedily by adding a particle of nitre. The blow-pipe is useful to the method is obtained as offering a ready method of knowing what the compourer parts of bodies are. Trials with the blow-pipe are generally made by the chemists in order to know the no rule of the constituent part, before to figure the need to the other steps of dry or humid analysis, which are requisite for ascertaining the quantities of the constituent rate. Then recourses is had to occur means than the blow-pipe; far in order to come as it a knowledge of the proportions of the constituent part, it is necessary that the quantity of each constituent part be large enough to be weighed in a balance, and, if in this purpose, the quantity of the substance employed must be larger than what can be manged with the blow-pipe. The first who applied the blow-pipe to the analysis of materials was weak, counselor of the college of mines in Sweden in 1725. Its application to the science of mineralogy was afterwards further improved by Cronstedt, Rine, may Gabon, Schede, and Bergman, and by other men of science since their time. In experimental intensity with the blow-pipe, the small fragment of the body subjected to not shell should its exceed the size of half a peppercorn; if larger, it cannot be sufficiently besoit. It is placed in a functional activity, made with a knife, in a piece of well-being charcoal to wood, free from gracks, and not too porous, and of the length of four or five inches, so as to be held conveniently Blow-pipe. In the left hand. Some blow-pipes have been made with a stand, to which they are connected by a ball and socket joint, also stands fixed to the table by a clampy and the construction leaves for light heat at liberty. In ycluding fragments of metallic ones by the blow-pipe, charcoal should be used as a support, for the charcoal attracts the oxygen from the metallic oxides, and reduces it to a metallic form; and when thus reduced, the metal may be kept fused on the charcoal, which prevents overrated its again attracting oxygen. The charcoal support has like› wise the advantage of increasing the heat by its hematdescence. For both these reasons, however, radiation, and to increase and reverberate the heat processing from the jet of flame, the goldishth wolders his small work by the blow-pipe attached to work to a piece of charcoal, by means of wires, in the process of soldering. When it is required that the frequency of a mineral Support of should be heated without the contact of charcoal, the plain. Engure is exposed to the flame in a small spoon of platina, with a wooden handle, the cavity of the spoon being a bremisphere of three inches of an inch in diameter; or on a thin lamina of plains, two or three inches long and half an inch broad, and of the thickness of common writing paper; or it is held by a foresee, three inches long, made of thin plasma. Bergens, who published his creation on the blow-pipe in 1768, before the working of platan had come into use, employed a small gold spoon, as that metal has the quality of remaining pine and uncontaminated, whilst in counter with many of the chemical agents; but platitan is preferable; for, besides possessing the quality of treating the surface of many changing areas, it has likewise the advantage of difficult fruiting. It has now for a good many years been brought into various instruments both in London and Paris, and why people it is sold at the price of about a guinea the ounce, which is one quarter the price of gold. Some platnary workers, as leancit of Paris, who was one of the firm for the crude gesimilar planatin into masses, by melting it with arsenic, and subsequent heating and/gong, where dissolve the crude plantin in interromulated soil, and reduce the nitrotranslate of platina to a metallic state by heat. Platina, however, although infilab alone by the heat of the common blowpipe, will be dissolved and melted if heated along with some of the metals. Platina supports, therefore, should not be used where they are liable to be in constant with a fused metal. These effects are notable in the case of tin; for when the is melted in contact with a vessel of platina, the the entire ironic combination with the platina, corroding and rendering it brittle, so that pieces of the platina vessel comes off on the application of a small force, and the vesad is thus rendered useless. Platina seeds also become unserivable by frequent and continued exposure to great heat. Platina crables which are much used because britthe, and carries it the diggers, and can should be taken to cool these vessels gradually, that they may last as long as possible. A platina vessel, which sulphites and was boiled for a long term, at least became performed and unserivecable. Borax (borate of soda) is used along with the fragment Fluxes of mineral in many cases. When exposed to the flame it become opaque, wells and ramides much, in consequence of parting with its water of crystallisation, increases it rises into a colorless and transparent lead. It is convenient to toe calcined borox, which is borax deprived of water or cyclic crystallization by heat in a crécyclical, this melts into a bead on the charcoal at once. The solubility of a natural in borx, with difference or without effervescents, and the colou which the mineral concentrates to the borax, are the chief distinctive characters obtained
\end{multicols}
\section*{Source record}
\begin{itemize}
\item Page ID: \texttt{britannica:309f765b1ecc83e0a7f7762e44b4a927a7555443:train:1818}
\item Scan: \texttt{page-1818-e6693d1961ef.jpg}, 600 x 816 px
\item Image SHA-256: \texttt{\small e6693d1961ef5657ea724a54f5c12cc8b62568f84a2ff587847275075a99e6a0}
\item Extraction: \texttt{nvidia/nemotron-3.5-parse}, raw reply \texttt{artifacts/extract/raw/1818.json}
\item Text is Parse OCR output, normalized (long s as s, line-end hyphens joined). It is not corrected by hand.
\end{itemize}
\end{document}
